\documentclass[12pt,a4paper]{article}
\usepackage[utf8]{inputenc}
\usepackage[T1]{fontenc}
\usepackage[ngerman,latin]{babel}
\usepackage{amsmath,amssymb,amsthm}
\usepackage{geometry}
\usepackage{hyperref}
\usepackage{enumitem}
\usepackage{mathrsfs}
\usepackage{longtable}
\usepackage{booktabs}
\usepackage{microtype}
\usepackage{array}
\usepackage{multirow}
\usepackage{graphicx}
\usepackage{colonequals}
\geometry{margin=2.5cm}

\theoremstyle{plain}
\newtheorem{satz}{Satz}[section]
\newtheorem{theorem}{Theorem}[section]
\newtheorem{lemma}{Lemma}[section]
\newtheorem{hilfssatz}{Hilfssatz}[section]
\newtheorem{korollar}{Korollar}[section]

\theoremstyle{definition}
\newtheorem{definition}{Definition}[section]

\theoremstyle{remark}
\newtheorem{bemerkung}{Bemerkung}[section]
\newtheorem{beispiel}{Beispiel}[section]

\newcommand{\origpage}[1]{\marginpar{\small\textit{S.~#1}}}

\title{Carl Friedrich Gau\ss\ Werke --- Band 11, Abt.~1, Teil 3}
\author{Carl Friedrich Gau\ss}
\date{}

\begin{document}
\maketitle

\section*{Physik. Briefwechsel und Nachlass}
\addcontentsline{toc}{section}{Physik. Briefwechsel und Nachlass}


% =====================================================
% [22.] GAUSS AN OLBERS UND GERLING
% =====================================================
\subsection*{[22.]}
\begin{center}
\textbf{Gau\ss\ an Olbers und Gerling.} G\"ottingen, 18.~April 1836.
\end{center}
\begin{center}
\rule{6cm}{0.4pt}\\[4pt]
[{\it Wilhelm Olbers, Sein Leben und seine Werke} II, 2, S.~636.]
\end{center}

\smallskip
\ldots\ldots\ldots\ldots\ldots\ldots\ldots\ldots\ldots\ldots\ldots\ldots
\smallskip

Sie theilen mir so g\"utig mit, da\ss\ Sie {\sc Weber}n in der neuen
Ausgabe seiner {\it Resultate} den Aufsatz: ``Neue Anwendungen des
Multiplicators'' erscheinen lassen. {\sc Weber} hat mir selbst davon
gesagt, aber, wie ich glaube, noch nichts davon an Sie abgeschickt. Er
hat viel bei diesen Sachen gearbeitet und die Resultate gehen weit
\"uber das hinaus, was in dem Programm enthalten ist. Er hat mehrere
hundert Vergleichungen der absoluten und durch den Multiplicator
erhaltenen Werthe gemacht. Ich habe die Sachen bei unserm
gegenw\"artigen Standpunkte der Magnetismus-Kenntnisse f\"ur sehr
wichtig gehalten und deshalb bei {\sc Weber} darauf gedrungen, die
Sache ordentlich und gr\"undlich herauszugeben, und nicht in
Extracten, wie er zuerst wollte. Ich habe ihm auch vorgeschlagen,
wie er die Formeln nach meiner Ansicht am zweckm\"a\ss igsten
einrichten soll. Ich hoffe, es wird eine sch\"one Abhandlung werden,
und werde die {\it Resultate} durch ein Vorwort einleiten.

\smallskip
\ldots\ldots\ldots\ldots\ldots\ldots\ldots\ldots\ldots\ldots\ldots\ldots
\smallskip

% =====================================================
% [23.] GAUSS AN SCHUMACHER
% =====================================================
\subsection*{[23.]}
\begin{center}
\textbf{Gau\ss\ an Schumacher.} G\"ottingen, 24.~Juni 1836.
\end{center}
\begin{center}
\rule{6cm}{0.4pt}\\[4pt]
[Briefwechsel zwischen Gau\ss\ und Schumacher III, S.~73.]
\end{center}

\smallskip
\ldots\ldots\ldots\ldots\ldots\ldots\ldots\ldots\ldots\ldots\ldots\ldots
\smallskip

{\sc Steinheil} f\"uhrt jetzt (als Probeversuch f\"ur k\"unftig
weiter zu erstreckende magnetische Telegraphie) eine Drahtleitung
von M\"unchen nach Bogenhausen, bei welcher Gelegenheit er schon
eine interessante Bemerkung gemacht hat --- meinem Vermuthen nach in
der atmosph\"arischen Electricit\"at begr\"undet, welche sich also an
den Magnetometern ausserordentlich stark sichtbar machen l\"asst.

\footnote{[*] Diese Zeitschrift erschien vom Juni 1837 an bis zum Jahre 1843 in 6
B\"anden unter der Ueberschrift: {\it Resultate aus den Beobachtungen des
magnetischen Vereins im Jahre} 1836, bezw.~1837, 1838, 1839, 1840, 1841.
Herausgegeben von C.~F.~Gau\ss\ und W.~Weber, G\"ottingen 1837, 1838,
Leipzig 1839, 1840, 1841, 1843.}

% =====================================================
% [24.] GAUSS AN OLBERS
% =====================================================
\subsection*{[24.]}
\begin{center}
\textbf{Gau\ss\ an Olbers.} G\"ottingen, 23.~Juli 1836.
\end{center}
\begin{center}
\rule{6cm}{0.4pt}\\[4pt]
[{\it Wilhelm Olbers, Sein Leben und seine Werke} II, 2, S.~642.]
\end{center}

\smallskip
\ldots\ldots\ldots\ldots\ldots\ldots\ldots\ldots\ldots\ldots\ldots\ldots
\smallskip

Unter meinen jetzigen magnetischen Experimenten ist das
merkw\"urdigste eines, wo die Inductionswirkung des Erdmagnetismus
auf einen $700$\,fachen Drahtring, \"ubersponnen und zusammen \"uber
$13000$ Fuss lang, bestimmt wird. Dieser Ring oder dieses Rad wird
um eine Axe gedrehet, die einen horizontalen Diameter des Rades
bildet. Diese Drehungsaxe macht genau einen rechten Winkel mit dem
magnetischen Meridian, und die beiden Enden des Drahts sind bis zum
Multiplicator des grossen Magnetometers der Sternwarte fortgef\"uhrt,
welcher Multiplicator jetzt aus $610$ Umwindingen \"ubersponnenen
Drahts besteht. Die Drehung geschieht tactm\"a\ss ig nach der Uhr,
alle $2$ Secunden Eine Umdrehung, und nach jeder halben Umdrehung
wechselt verm\"oge eines eigenth\"umlichen Mechanismus die
Verbindung der Drahtendungen des Rades mit ihren Fortsetzungen. Die
Einrichtung ist nach Vorschrift der Theorie so, da\ss\ dieser Wechsel
genau Statt findet, wenn das Rad dem magnetischen Aequator parallel
ist, oder einfacher gesagt, normal gegen die Richtung der
Erdmagnetischen Kraft. So wirkt der durch die Induction erzeugte
Strom immer in Einerlei Sinn auf das Magnetometer und bewirkt,
obgleich die ganze Drahtl\"ange so gegen $20000$ Fuss lang ist, doch
noch Ausweichung von mehreren hundert Scalentheilen an der
$25$\,pf\"undigen Nadel. Bei einigen Versuchen wird auch noch der
Draht des Inductors ({\sc Schumachers} J.~B.~1836, S.~41\footnote{[*]
{\it Erdmagnetismus und Magnetometer}, Jahrbuch f\"ur 1836, herausgeg.~von
Schumacher, Werke V, S.~315, siehe insbesondere S.~340.}), der aber
jetzt nicht mehr $3537$, sondern $7000$ Umwindingen hat, mit in die
Kette aufgenommen, wo also der Strom gegen $27000$ pariser Fuss sehr
d\"unnen Drahts zu durchlaufen hat und dann noch immer $325$
Scalentheile Ausschlag gibt. Diese Versuche machen einen Theil von
denjenigen aus, die zum Zweck haben, alles was sich auf die
Wechselwirkung zwischen Magnetismus und galvanischen Str\"omen
bezieht, auf absolute Maasse zu bringen. Gen\"aherte
Zahlenbestimmungen geben schon meine fr\"uheren Versuche. Die
gegenw\"artigen werden aber die Genauigkeit noch sehr vergr\"ossern.
\ldots\ldots\ldots\ldots\ldots\ldots\ldots\ldots\ldots\ldots\ldots\ldots

% =====================================================
% [25.] GAUSS AN GERLING
% =====================================================
\subsection*{[25.]}
\begin{center}
\textbf{Gau\ss\ an Gerling.} G\"ottingen, 19.~December 1836.
\end{center}

\smallskip
\ldots\ldots\ldots\ldots\ldots\ldots\ldots\ldots\ldots\ldots\ldots\ldots
gewesen. \ldots\ldots\ldots\ldots\ldots\ldots\ldots\ldots\ldots\ldots

Mein Sohn ist jetzt wieder in Stade; er hat sich 3 Wochen in Paris
aufgehalten und hat von den dortigen Gelehrten viel Gef\"alligkeit
erfahren, besonders von {\sc Poisson} und {\sc Libri}. Auch mit
{\sc Arago} ist er bekannt geworden, welcher ihm erkl\"art hat, er
ziehe doch die {\sc Gambey}schen Apparate den meinigen vor, weil
letztere zu viel Anomalien zeigten. {\sc Airy}'s Bedenklichkeiten
bestanden hingegen, wie Sie sich erinnern, darin, da\ss\ sie meinten,
meine Apparate zeigten zu wenig Anomalien.
\ldots\ldots\ldots\ldots\ldots\ldots\ldots\ldots\ldots\ldots\ldots\ldots

% =====================================================
% [26.] GAUSS AN OLBERS
% =====================================================
\subsection*{[26.]}
\begin{center}
\textbf{Gau\ss\ an Olbers.} G\"ottingen, 2.~September 1837.
\end{center}
\begin{center}
\rule{6cm}{0.4pt}\\[4pt]
[{\it Wilhelm Olbers, Sein Leben und seine Werke} II, 2, S.~649.]
\end{center}

\smallskip
\ldots\ldots\ldots\ldots\ldots\ldots\ldots\ldots\ldots\ldots\ldots\ldots
gro\ss en Werth lege. \ldots\ldots\ldots\ldots\ldots\ldots\ldots

In der Sitzung der Societ\"at im Jubil\"aum (19.~September) werde ich
eine Vorlesung halten \"uber ein neues Mittel f\"ur die Magnetischen
Beobachtungen\footnote{[*] Siehe den Bericht \"uber die am 19.~September 1837 gehaltene
Vorlesung: {\it Ein neues H\"ulfsmittel f\"ur die Magnetischen
Beobachtungen} in den G\"ottingischen Gelehrten Anzeigen vom
30.~Oktober 1837, Werke V, S.~352, und den Aufsatz {\it \"Uber ein
neues, zun\"achst zur unmittelbaren Beobachtung der Ver\"anderungen
in der Intensit\"at des horizontalen Theils des Erdmagnetismus
bestimmtes Instrument} in den Resultaten aus den Beobachtungen des
magnetischen Vereins im Jahre 1837, S.~1, Werke V, S.~357.}. Es
bezieht sich auf einen neuen Apparat, der f\"ur die (horizontale)
Intensit\"at ganz dasselbe leistet, was das Magnetometer f\"ur die
Declination, wodurch also die Aufgabe ({\it Resultate}, S.~12\footnote{[**]
{\it Einleitung}, Resultate aus den Beobachtungen des magnetischen Vereins
im Jahre 1836, S.~1, Werke V, S.~345, siehe insbesondere S.~351.}),
soweit von dem horizontalen Theile der Erdmagnetischen Kraft die
Rede ist, erledigt wird.

Von einer rohen Probe der Grundidee finden Sie eine Andeutung in
{\sc Schumachers} Jahrbuch 1836, S.~19\footnote{[***] {\it
Erdmagnetismus und Magnetometer}, Werke V, S.~315, siehe insbesondere
S.~326 Zeile 4 v.~u.~bis S.~327 Zeile 2.}, woraus freilich nicht zu
erkennen ist, in was das Mittel besteht, sondern nur ein secund\"arer
Theil von dem, was damit geschieht. In diesem Sommer habe ich aber
den Apparat ordentlich ausf\"uhren lassen, und es sind sogar die
beiden letzten Termine (der Haupttermin vom 29.~Juli und der
Extrater\-min vom 31.~August) vollst\"andig damit beobachtet,
w\"ahrend ebenso vollst\"andig in beiden auch im M[agnetischen]
O[bservatorium] der Verlauf der Declination beobachtet ist.

Der horizontale Theil des Erdmagnetismus kann also jetzt so scharf
beobachtet werden, wie die Sterne am Himmel. Aber mit dem verticalen
Theile wird eine \"ahnliche Genauigkeit niemals erreicht werden k\"onnen;
wer die Stelle der {\it Intensitas vis etc.}~p.~15 ``{\it Ex hoc rem
--- requiratur}''\footnote{[***] Werke V, S.~91 letzter Absatz.}
geh\"orig studirt und beherzigt hat, wird dies leicht von selbst
einsehen. Wenn man aber auch nicht die gleiche Genauigkeit wie bei
dem horizontalen Theil erreichen kann, so bin ich doch \"uberzeugt,
da\ss\ man viel mehr erreichen kann, als bis jetzt erreicht ist. Das
ist aber etwas, worauf ich mich nicht einlassen kann. Bei den
Instrumental-H\"ulfsmitteln f\"ur die beiden Elemente des horizontalen
Theils konnten geistige Mittel ausreichen, d.~i.~durch eine geh\"orige
Einrichtung konnte man Apparate zur Erreichung der h\"ochsten
Genauigkeit darstellen, die eigentlich keine \"uberm\"a\ss ig feine
mechanische Arbeit erfordern, und mit geringen Kosten angefertigt
werden k\"onnen. (Der neue Apparat kommt nicht auf 50\,\textdollaroldstyle,
nat\"urlich alles schon vorhandene, was dabei gebraucht wird,
ungerechnet, namentlich Uhr, Theodolith und Magnetstab.) Dagegen
sind bei allem, wobei der verticale Theil der magnetischen Kraft auf
irgend eine Art mit ins Spiel kommt, sei es als Inclination oder
anders, sehr fein, sehr vollkommen ausgearbeitete, also auch am Ende
kostbare Instrumente unentbehrlich, unerl\"a\ss lich und k\"onnen durch
Nichts anderes ersetzt werden. Ein Inclinationsapparat von der
gew\"ohnlichen Einrichtung, der nur so weit befriedigen soll, als wir
jetzt wirklich sind, also wobei man noch weit, sehr weit zur\"uck ist
gegen das, was man w\"unschen mu\ss, kostet schon mehr, als ich in
einem oder zwei Jahren f\"ur Sternwarte und Magnetische Anlagen
zusammen zu verausgaben habe.
\ldots\ldots\ldots\ldots\ldots\ldots\ldots\ldots\ldots\ldots\ldots\ldots


% =====================================================
% [27.] GAUSS AN OLBERS
% =====================================================
\subsection*{[27.]}
\begin{center}
\textbf{Gau\ss\ an Olbers.} G\"ottingen, 20.~November 1838.
\end{center}
\begin{center}
\rule{6cm}{0.4pt}\\[4pt]
[{\it Wilhelm Olbers, Sein Leben und seine Werke} II, 2, S.~697.]
\end{center}

\smallskip
\ldots\ldots\ldots\ldots\ldots\ldots\ldots\ldots\ldots\ldots\ldots\ldots
\smallskip

Ich gehe jetzt damit um, den Apparat, dessen in den ``Resultaten
f\"ur 1837'' p.~7--8\footnote{[*] Siehe den oben S.~111 in der Fu\ss note
genannten Aufsatz, insbesondere Werke V, S.~362.} erw\"ahnt ist, auf
eine eigenth\"umliche Art ausf\"uhren zu lassen. Die Ab\"anderung,
welche {\sc Weber} unter dem Namen Inductions-Inclinatorium gemacht
hat, ist zwar im h\"ochsten Grade sinnreich; ich glaube aber nicht,
da\ss\ man auf diese Art jemals sehr scharfe Resultate erhalten kann.
Ich komme auf meine alte Art zur\"uck (wobei ein von dem
Drehungsapparat ganz getrenntes Magnetometer gebraucht wird), lasse
aber an jenem zwei getheilte Kreise und Vorkehrungen zum scharfen
Nivelliren anbringen. Ich bin geneigt zu glauben, da\ss\ man damit
ziemlich scharfe Inclinationen, wenigstens so scharfe, wie mit den
gew\"ohnlichen Inclinatorien, erhalten kann, ja, wenn man noch anderes
damit verbinden will, auch absolute Declination. Doch werde ich erst
den Erfolg abwarten, ehe ich mich weiter dar\"uber \"aussere.
\ldots\ldots\ldots\ldots\ldots\ldots\ldots\ldots\ldots\ldots\ldots\ldots

% =====================================================
% [28.] GAUSS AN ENCKE
% =====================================================
\subsection*{[28.]}
\begin{center}
\textbf{Gau\ss\ an Encke.} G\"ottingen, 8.~September 1839.
\end{center}

\smallskip
\ldots\ldots\ldots\ldots\ldots\ldots\ldots\ldots\ldots\ldots\ldots\ldots
Arbeit ein. \ldots\ldots\ldots\ldots\ldots\ldots\ldots\ldots\ldots

Das g\"utige Interesse, welches Sie an meiner allgemeinen Theorie des
Erdmagnetismus bezeigen, ist mir sehr erfreulich gewesen. Wenn diese
Arbeit ein Verdienst hat, so besteht es meiner Meinung nach darin,
da\ss\ man fr\"uher nicht wu\ss te, wie man \"uberhaupt die Sache
anzufassen habe, und da\ss\ jetzt der Weg dazu offen ist. Ich habe
insofern eine numerische Anwendung ungern gemacht, als ich die
Ueberzeugung hatte und habe, da\ss\ die groben Data jetzt eine so
gro\ss e mechanische Arbeit, wie ich habe aufwenden m\"ussen, kaum
verdienen, und da\ss\ bei guten Datis eine nicht l\"angere Arbeit
ganz anders befriedigende Resultate h\"atte liefern k\"onnen und
m\"ussen. Indessen war ich eben so sehr \"uberzeugt, da\ss\ eine
blosse Auseinandersetzung der Theorie gar keinen Eingang gefunden
haben w\"urde.

In Beziehung auf Ein Moment scheinen Sie \"ubrigens sich nicht ganz
in den eigentlichen Standpunkt gesetzt zu haben, n\"amlich da\ss\ Sie
pag.~47\footnote{[*] Es handelt sich genauer um Seite 46 und 47 der
Resultate aus den Beobachtungen des Magnetischen Vereins im Jahre 1838, wo
im Artikel 32 der {\it Allgemeinen Theorie des Erdmagnetismus} (siehe
Werke V, S.~162) das Theorem des art.~2. der {\it Intensitas vis
magneticae terrestris} erw\"ahnt wird, vergl. oben S.~103 die zweite
Fu\ss note.} einen Anstoss gefunden haben. Sie scheinen dieses nur
wie einen kleinen Nebenpunkt zu betrachten, wo Ihnen nicht gleich
gelungen sei die L\"ucke zu suppliren. Es ist aber gar nicht meine
Meinung gewesen, da\ss\ die Sache aus diesem Gesichtspunkt betrachtet
werden soll. Vielmehr ist diess nur Ein Hauptsatz aus einer
gr\"osseren theoretischen Untersuchung, auf die ich selbst nach meinen
eigenen Maassstabe, einen viel gr\"osseren Werth lege, als auf die
ganze Abhandlung, die jetzt gedruckt ist, insofern in jener Arbeit
$10$~mahl mehr geistige von einander unabh\"angige Schritte zu machen
sind, als in dieser, wo eigentlich nach der einmahl gefassten
Grund-Idee alles sich von selbst macht.
\ldots\ldots\ldots\ldots\ldots\ldots\ldots\ldots\ldots\ldots\ldots\ldots

% =====================================================
% [29.] GAUSS AN HERGER
% =====================================================
\subsection*{[29.]}
\begin{center}
\textbf{Gau\ss\ an Herger.}
\end{center}
\begin{center}
Wohlgeborener\\
Hochzuehrender Herr.
\end{center}

F\"ur das g\"utige Geschenk, welches Sie mir mit der ersten Lieferung
Ihres sch\"atzbaren auch in seinem Aeusseren so elegant
ausgestatteten Werkes gemacht haben\footnote{[*] J.~{\sc Ernst Herger},
{\it Die Systeme der magnetischen Curven, Isogonen und Isodynamen,
nebst anderweitigen empirischen Fortschritten \"uber die
magnetisch-polaren Kr\"afte, ausgef\"uhrt in 37 grossen graphischen
Darstellungen auf 31 Tafeln und erl\"autert unter den Auspicien des
Hofrath Dr.~Schottin}. Nebst einem Vorwort von G.~A.~{\sc Erman}.
In 4 Heften, Leipzig 1845. --- Der Verfasser dieses Werkes
J.~E.~{\sc Herger} (1812--1871) lebte als Rosenz\"uchter und
Handelsg\"artner zu K\"ostritz im F\"urstentum Reu\ss.},
verfehle ich nicht, Ihnen meinen schuldigen verbindlichsten Dank
abzustatten. Der von Ihnen dabei bewiesene ausdauernde Fleiss, sowie
die Accuratesse und Nettigkeit der Ausf\"uhrung gereichen Ihnen zur
Ehre und verdienen alle Anerkennung. Ihre Curven sind in der That
nichts anderes als die rechtwinklichten Trajectorien (denn so wird
ohne Zweifel das unleserlich geschriebene Wort in Herrn von
{\sc Humboldts} Briefe gelesen werden m\"ussen) zu denjenigen Curven,
welche nach meiner Theorie gleiche Potentiale vorstellen. Aber die
Gestaltung des einen wie des andern Curvensystems in der N\"ahe von
k\"unstlichen Magneten l\"a\ss t sich theoretisch oder {\it a priori}
nicht bestimmen, ohne die Vertheilung des Magnetismus im Innern von
diesen vorher zu kennen, und so bieten, wie Hr.~Pr[ofessor]~{\sc Ermann}
richtig bemerkt, Ihre Curven ein sehr sch\"atzbares Material dar,
wonach jede vorgebrachte oder vorzubringende Hypothese \"uber solche
Vertheilung wird gepr\"uft werden k\"onnen.

D\"urfte ich mir erlauben, noch einen Wunsch auszusprechen, so w\"are
es der, da\ss\ in der Einleitung \"uber die Versuche selbst, auf welche
die gezeichneten Curven sich gr\"unden, ein vollst\"andigeres Detail
mitgetheilt w\"urde, ein Wunsch, den Sie vielleicht schon selbst im
Fortgange der Einleitung zu erf\"ullen sich vorgesetzt haben, da dies
in der That das am meisten geeignete Mittel ist, das Zutrauen zu
Ihren Resultaten zu befestigen und zu erh\"ohen.

Genehmigen Sie die Bezeugung der besonderen Hochachtung womit ich
beharre Ewr. Wohlgeboren

\begin{flushright}
ergebenster Diener\\
\textsc{C. F. Gau\ss.}
\end{flushright}
\begin{center}
G\"ottingen, 11. Februar 1845.
\end{center}

% =====================================================
% [30.] GAUSS AN BESSEL
% =====================================================
\subsection*{[30.]}
\begin{center}
\textbf{Gau\ss\ an Bessel.} G\"ottingen, 31.~December 1831.
\end{center}
\begin{center}
\rule{6cm}{0.4pt}\\[4pt]
[Briefwechsel zwischen Gau\ss\ und Bessel, Berlin 1880, S.~504.]
\end{center}

\smallskip
\ldots\ldots\ldots\ldots\ldots\ldots\ldots\ldots\ldots\ldots\ldots\ldots
\smallskip

Bei der Lect\"ure von {\sc Poissons} \"alterer Abhandlung \"uber die
Electricit\"at (1812)\footnote{[*] S.~D.~{\sc Poisson}, {\it M\'emoire
sur la distribution de l'Electricit\'e \`a la surface des corps
conducteurs}, M\'emoires de l'Institut, ann\'ee 1811, Paris 1812, S.~1.}
bin ich neulich auf einen meines Wissens neuen und wie mir scheint
artigen Satz gekommen. Er geht [a.~a.~O.~S.~3] davon aus, da\ss\ die
freie Electricit\"at sich unter der Oberfl\"ache eines Isolirten
Leiters sammelt, und meint, es sei nicht zureichend, da\ss\ die
Resultante aller Abstossungen in jedem Punkte der inneren
Oberfl\"ache jener Electricit\"atsschicht auf diese senkrecht sei,
sondern die physische Aufgabe fordere ausserdem noch die Bedingung,
da\ss\ die Resultante in jedem Punkte des innern eingeschlossenen
Raums $=0$ werde, weil sich sonst neue Electricit\"at zersetzen
w\"urde. Man mag diese Behauptung an sich zugeben, allein ich
{\it beweise}, da\ss, unter Voraussetzung der Repulsion im
verkehrten doppelten Verh\"altni\ss\ des Abstandes, die zweite
Bedingung schon von selbst in der ersten enthalten ist. Dann wird
es offenbar ein Fehler, da von neuem an die physische Natur zu
appelliren, wo die Mathematik den Gegenstand schon von selbst
darreicht. Mein Satz gilt eben so gut von anziehenden Kr\"aften und
heisst dann: Wenn ein wie immer gestalteter homogener oder nicht
homogener K\"orper aus Theilen besteht, die eine Anziehungskraft im
verkehrten doppelten Verh\"altni\ss\ des Abstandes aus\"uben, und
dieser K\"orper einen hohlen Raum umschliesst, an dessen Begrenzung
in keinem Punkte eine schiefe Resultante Statt hat, so ist die
Resultante an jedem Punkte dieser Begrenzung sowohl als in jedem
Punkte des hohlen Raums $=0$. --- Sobald man den Satz einmahl
aufgefasst hat, ist die Beweisf\"uhrung eben nicht schwer zu finden.
\ldots\ldots\ldots\ldots\ldots\ldots\ldots\ldots\ldots\ldots\ldots\ldots


\newpage
\section*{[Dioptrik.] Nachlass.}
\addcontentsline{toc}{section}{[Dioptrik.] Nachlass.}

% =====================================================
% [I.] NEUER ALGORITHMUS
% =====================================================
\subsection*{[I.]}
\begin{center}
\textbf{NEUER ALGORITHMUS.}
\end{center}
\begin{center}
\rule{6cm}{0.4pt}\\[4pt]
[Aus Handbuch 21, Bg, Aufs\"atze, Notizen und Rechnungen zur
Mathematik und Astronomie geh\"orig; Angefangen September 1813, S.~27,
28.]
\end{center}

Die Gr\"ossen $A$, $B$, $C$, $D$ etc.~entstehen aus $a$, $b$, $c$, $d$
etc.~so, da\ss
\[
A = a,\quad B = bA - 1,\quad C = cB - A,\quad D = dC - B\ \text{ etc.}
\]
Wir bezeichnen dann
\[
A\ \text{mit}\ (a,a),\quad B\ \text{mit}\ (a,b),\quad C\ \text{mit}\ (a,c),\quad
D\ \text{mit}\ (a,d)\ \text{etc.},
\]
wo also bei $a$, $b$, $c$, $d$ etc.~die Ordnung als bestimmt
betrachtet wird.

Dieser Algorithmus bietet die interessante Eigenschaft, da\ss\ man\footnote{[*]
In der Handschrift steht hier noch das Wort ``statt''.} vier oder
mehr aufeinander folgende Elemente immer durch drei ersetzen kann.

Man hat dann folgende sch\"one Lehrs\"atze (Ordn[un]g $a$, $b$, $c$,
\ldots $g$, $h$, $i$, \ldots $p$, $q$, $r$, \ldots $z$)
\begin{align*}
\text{I.}\qquad (a,p)\cdot(i,z) - (i,p)\cdot(a,z) &= (a,g)\cdot(r,z),\\
\text{II.}\qquad \partial\,(a,z) &= (b,z)\,\partial a + (a,a)(c,z)\,\partial b
+ (a,b)(d,z)\,\partial c + (a,c)(e,z)\,\partial d + \text{etc.}
\end{align*}

\bigskip

Diese Formeln werden noch eleganter, wenn man die Bezeichnung\footnote{[*]
In der Handschrift steht hier noch das Wort ``anders''.} modifieirt und
\[
(a,a) = 0,\quad (a,b) = 1,\quad (a,c) = b,\quad
(a,d) = c\,(a,c) - (a,b),\quad (a,e) = d\,(a,d) - (a,c)\ \text{etc.}
\]
setzt. \textbf{Alsdann hat man nemlich:}
\begin{align*}
\text{I[a].}\qquad (a,p)\cdot(i,z) - (i,p)\cdot(a,z) &= (a,i)\cdot(p,z),\\
\text{II[a].}\qquad \partial\,(a,z) &=
(a,b)(b,z)\,\partial b + (a,c)(c,z)\,\partial c + (a,d)(d,z)\,\partial d
+ \text{etc.}
\end{align*}

Auf Einer Axe seyen $\mu + 1$ Gl\"aser geordnet, deren Brennweiten
$\frac{1}{f}$, $\frac{1}{f'}$, $\frac{1}{f''}$,
$\frac{1}{f'''}$ etc.~$\frac{1}{f^{\mu}}$, die Distanzen von
einander $h$, $h'$, $h''$, \ldots $h^{\mu-1}$ sind. Man mache,
indem obigen Elementen $a$, $b$, $c$ etc.~jetzt die Gr\"ossen $f$,
$h$, $f'$, $h'$, $f''$, $h''$ etc.~entsprechen:
\[
\begin{array}{ll}
\text{[III.]}\ \left\{\begin{array}{l}
\alpha = 1,\\
\alpha' = (f,h),\\
\alpha'' = (f,h'),\\
\alpha''' = (f,h''),\end{array}\right.\qquad
&
\text{[IV.]}\ \left\{\begin{array}{l}
\beta = f,\\
\beta' = (f,f'),\\
\beta'' = (f,f''),\\
\beta''' = (f,f'''),\\
\text{etc.}\end{array}\right.\\[2ex]
\text{[V.]}\ \left\{\begin{array}{l}
\gamma = 0,\\
\gamma' = h,\\
\gamma'' = (h,h'),\\
\gamma''' = (h,h''),\end{array}\right.\qquad
&
\text{[VI.]}\ \left\{\begin{array}{l}
\delta = 1,\\
\delta' = (h,f'),\\
\delta'' = (h,f''),\\
\delta''' = (h,f'''),\\
\text{etc.}\end{array}\right.
\end{array}
\]

\textbf{Die Bedeutung dieser Gr\"ossen versinnlicht folgende Figur:}

\medskip

[Fig.~A und B werden durch die in den Original angegebenen Schemata
dargestellt. In Fig.~A hat der Strahl die Richtung 1, 2, 3, \ldots;
in Fig.~B die Richtung 1, 2, 3, \ldots, und die\vspace{-3pt}

$\alpha$, $\alpha'$, $\alpha''$ etc.\ bezeichnen die St\"ucke des Weges
vom Punkt 1 bis zu 2, von 2 bis 3 etc.,

$\beta$, $\beta'$, $\beta''$ etc.\ die Verh\"altnisse von $y^{(1)}$ zu $y$,
$y^{(2)}$ zu $y$ etc.,

$\gamma$, $\gamma'$, $\gamma''$ etc.\ die St\"ucke von der gemeinschaftlichen
Axe der Linsen ab gerechnet,

$\delta$, $\delta'$, $\delta''$ etc.\ die Verh\"altnisse von $\eta^{(1)}$ zu $\eta$,
$\eta^{(2)}$ zu $\eta$ etc.]

\medskip

Hier wird nun:

\quad $\delta^{\mu}$ die Vergr\"osserung; bei geradem $\mu$ ist
$\delta^{\mu}$ positiv bei aufrechtem Bilde, negativ bei verkehrtem,

\quad $\frac{g}{\beta}$, $\frac{g}{\beta'}$, $\frac{g}{\beta''}$ etc.\qquad
Gr\"osse der einzelnen Bilder, wenn $g$ das Gesichtsfeld,

\quad $g\gamma'$, $g\gamma''$, $g\gamma'''$ etc.\qquad
Oeffnungen der Gl\"aser wegen des Gesichtsfeldes,

\quad $\frac{\omega}{\delta}$, $\frac{\omega}{\delta'}$,
$\frac{\omega}{\delta''}$ etc.\qquad
Durchmesser der Strahlenkegel,

\quad $\omega\alpha'$, $\omega\alpha''$, $\omega\alpha'''$ etc.\qquad
Oeffnungen der Gl\"aser wegen der Helligkeit.

\medskip

F\"ur weitsichtige Augen mu\ss\ $\beta^{\mu} = 0$ werden; f\"ur
kurzsichtige, deren Augenweite $= k$, mu\ss\ $\beta^{\mu}\delta^{\mu}k = -1$
seyn.

Die Breite des farbigen Saumes um das Gesichtsfeld wird
\[
\text{[VII]}\qquad \tfrac{1}{2}\delta^{\mu}g\left(\alpha^{\mu}\,\partial\delta^{\mu}
- \beta^{\mu}\,\partial\gamma^{\mu}\right)
\]
und zwar ist der Saum, wenn diese Gr\"os\ss e positiv wird, blau; roth,
wenn sie negativ ist\footnote{In der Handschrift steht in den drei
letzten Formeln $\delta^{m}$ statt $\delta^{\mu}$.}
\begin{align*}
\text{[VIIa]}\qquad &= \tfrac{1}{2}\delta^{\mu}g\,\frac{\partial n}{n-1}
\left\{\alpha'\gamma'f' + \alpha''\gamma''f'' + \alpha'''\gamma'''f'''
+ \alpha^{\text{IV}}\gamma^{\text{IV}}f^{\text{IV}}
+ \text{etc.} + \alpha^{\mu}\gamma^{\mu}f^{\mu}\right\},\\
\text{[VIIIb]}\qquad &= \tfrac{1}{2}\delta^{\mu}g\,\frac{\partial n}{n-1}
\left\{\alpha'(\delta + \delta') + \alpha''(\delta' + \delta'')
+ \alpha'''(\delta'' + \delta''') + \text{etc.}
+ \alpha^{\mu}(\delta^{\mu-1} + \delta^{\mu})\right\},\\
\text{VIIc]}\qquad &= \tfrac{1}{2}\delta^{\mu}g\,\frac{\partial n}{n-1}
\left\{\gamma'(\beta + \beta') + \gamma''(\beta' + \beta'')
+ \gamma'''(\beta'' + \beta''') + \text{etc.}
+ \gamma^{\mu}(\beta^{\mu-1} + \beta^{\mu})\right\}.
\end{align*}


% =====================================================
% [II.] ANALYSE VON FERNROEHREN
% =====================================================
\subsection*{[II.]}
\begin{center}
\textbf{ANALYSE VON FERNROEHREN.}
\end{center}
\begin{center}
\rule{6cm}{0.4pt}\\[4pt]
[Aus Handbuch 19, das "ubrige von dem Aufsatze Anfang April 1811 bis
Ende September 1814; S.~38--41.]
\end{center}

Bei der Aufsuchung verschiedener Fernr\"ohre kann die Formel f\"ur die
Feinheit der Theilung auf folgende Art umgeformt werden:
\[
\frac{g'}{\beta\beta'\cdot\sqrt{\alpha'}} =
\frac{g''}{\beta'\beta''\cdot\sqrt{\alpha''}} =
\frac{g'''}{\beta''\beta'''\cdot\sqrt{\alpha'''}} = \text{etc.}
\]
Hier bedeuten die Gr\"ossen $\alpha$, $\alpha'$, $\alpha''$, $\alpha'''$
etc.~die Distanzen der Linsen\footnote{In der Handschrift fehlt das
Wort ``Linsen''.} von einander, gemessen in solchen Einheiten, in
welchen der Abstand des betreffenden Bildes von dem entgegengesetzten
Endpunkte\footnote{In der Handschrift steht statt ``entgegengesetzten
Endpunkte'' wohl irrt\"umlich ``Endpunkte der Schwerpunkte''.} des
betreffenden Teleskop-Theiles als Einheit angenommen wird.

F\"ur die Unsch\"arfe hat man folgende Formel:
\[
\tfrac{1}{4}\delta^{\mu}g\cdot\partial\frac{\beta^{\mu}}{\delta^{\mu}}
\]
oder
\[
\tfrac{1}{4}g\left(\beta^{\mu}\partial\delta^{\mu}
- \delta^{\mu}\partial\beta^{\mu}\right).
\]

\subsubsection*{Beobachtungen.}
\begin{enumerate}
\item Wenn die Distanzen der Linsen als gegeben und f\"ur das Objektiv
auch die Brennweite $= f$, f\"ur die \oe ibrigen die letzteren und die
\oe ibrigen Verh\"altnisse, als gesucht betrachtet werden: so werden in
der Regel drei Linsen das H\"ochste leisten. F\"ur $y^{(1)}$, $y^{(2)}$,
$y^{(3)}$ werden die Formeln sehr einfach, n\"amlich
\[
y^{(1)} = y^{(2)} = y^{(3)} = \sqrt[3]{y^{2}f}.
\]

\item Wenn dagegen sowohl die Distanzen als auch die Brennweiten der
Linsen, mit Ausnahme der \oe ussersten beiden, sowohl ganz als
theilweise f\"ur das \oe usserste Augengl\oe as, gegeben sind, und
alles \oe ibrige gesucht wird: so wird f\"ur die \oe usserste Linse,
welche ein biconvexes Objektiv repr\"asentiert, zu dem Augengl\oe as
eine einfache divergirende Linse der Brennweite $f'$ eingeschaltet, und
es wird
\[
y^{(2)} = \tfrac{1}{2}y\sqrt[3]{\frac{y^{(1)2}}{y^{2}f}},\qquad
h' = \tfrac{3}{2}y^{(2)},\qquad
y^{(3)} = y^{(1)}\sqrt[3]{\frac{y^{2}f}{y^{(1)2}}},
\]
w\"ahrend das Gesichtsfeld $= 1\frac{1}{2}\,^{\prime}$ betr\"agt.
\end{enumerate}

\textbf{Zum \S 3.} \ Anwendung des letzteren Verfahrens auf ein
anastigmatisches Fernrohr mit einem einfachen und einem dreifachen
Objektive: Dabei k\"onnen drei und mehr Linsen mit Vorteil verwendet
werden, wobei man f\"ur achromatische Teleskope Formeln ableitet, die
die Positionen der Brennpunkte, die Vergr\"osserungen und die
chromatischen Bedingungen genau bestimmen.


% =====================================================
% [III.] DIOPTRIK
% =====================================================
\subsection*{[III.]}
\begin{center}
\textbf{DIOPTRIK.}
\end{center}
\begin{center}
\rule{6cm}{0.4pt}\\[4pt]
[Aus Handbuch 21, als "uberschrift des Blattes 14; das "ubrige S.~82.]
\end{center}

\noindent
(\textit{Bei einem dreyfach-objektivischen Fernrohre mit \oe ffner,
dem ein einfacher biconvexer Linsenk\"orper, ein zweifach achromatisches
Objektiv und ein dreifach achromatisches Fernrohr vorangeht.})

\medskip

\noindent
\textbf{1.} Es werde die Entfernung der beiden \oe ussersten Linsen
$= v + t$, oder vielmehr $= v + vt$, gesetzt, wo $v$ eine sehr kleine
Gr\"o\ss e gegen die Brennweite ist.

\noindent
\textbf{2.} Die Differenz von der Brennweite wird dadurch bestimmt,
da\ss\ man setzt
\[
\frac{1}{f^{(1)}} = \frac{1}{v} + \frac{1}{t} + 1.
\]

\noindent
\textbf{3.} Sodann wird die L\"osung der Gleichung f\"ur die
Verh\"altnisse von $v$ und $t$ gegeben.

\medskip

\noindent
\textbf{Zum \S 2.} \ Bei achromatischen Objektiven werden die
Formeln f\"ur die Brennweite aus den Brechungsverh\"altnissen abgeleitet.

\noindent
\textbf{Zum \S 4.} \ F\"ur die zusammengesetzten Objective kann man
"ahnliche Formeln aufstellen. [Die Fig.~E.~F.~G.~H.~J.~K.~beziehen
sich auf die verschiedenen Glassorten.]

\noindent
\textbf{Zum \S 5.} \ F\"ur Kronglas $n = 1.53$, f\"ur Flintglas
$n = 1.58$.

\noindent
\textbf{Zum \S 6.} \ Formeln f\"ur achromatische Fernr\"ohre:
\[
\frac{1}{f'} = \frac{1}{f\vphantom{()}} - \frac{1}{f''},
\]
wobei die Distanz $d$ der beiden Linsen so gew\"ahlt wird, da\ss\ die
chromatische Aberration verschwindet.

\noindent
\textbf{Zum \S 7.} \ Bei einem dreiachrigen Fernrohre sind die
Bedingungen so zu stellen, da\ss\ sowohl die erste als die zweite
Differenz der Brennweiten verschwindet.

\noindent
\textbf{Zum \S 8.} \ F\"ur eine ganz aplanatische und achromatische
Linse (Fig.~L) ist zu verlangen, da\ss\ sowohl die sph\"arische
Aberration, als die chromatische verschwinden. Zwei solche Gl\"aser
werden also 4 Gleichungen geben. Aber die L\"osungen k\"onnen entweder
unm\"oglich werden, oder Linsen geben, die nicht m\"oglich sind, und
allein die Glassorte ist willk\"urlich.

\noindent
\textbf{Zum \S 9.} \ N\"amlich setzt man die Brennweite des ersten
Glases $= b$ und des zweiten $= b'$, so da\ss\ $b - b'$ die Differenz
der reciproken Brechungsverh\"altnisse gibt, und bezeichnet durch $f$,
$f'$ die Oeffnungen, durch $r$, $r'$ die Radien, so ergeben sich
folgende Gleichungen:
\[
\frac{1}{b} = (n-1)\left(\frac{1}{r} - \frac{1}{r'}\right),\qquad
\frac{1}{b'} = (n'-1)\left(\frac{1}{r} - \frac{1}{r'}\right).
\]

\noindent
\textbf{Zum \S 10.} \ In jeder Reihe von Fernr\"ohren wird f\"ur ein
gew\"ohnliches Fernrohr die gr\"o\ss te und f\"ur ein aplanatisches die
kleinste Oeffnung erhalten.


% =====================================================
% [IV.] AUFGABE AUS DER DIOPTRIK
% =====================================================
\subsection*{[IV.]}
\begin{center}
\textbf{AUFGABE AUS DER DIOPTRIK.}
\end{center}
\begin{center}
\rule{6cm}{0.4pt}\\[4pt]
[Aus Handbuch 21, das "ubrige von dem Aufsatze, Anfang April 1811
bis Ende September 1814; S.~41.]
\end{center}

\textit{Zwei Gl\"aser sind gegeben: man soll diejenige Linse finden,
welche mit jenen verbunden ein Fernrohr von gegebener L\"ange giebt,
und wobei der Farbenfehler ein Minimum wird.}

Es seien die Brechungsverh\"altnisse der beiden gegebenen Gl\"aser
$n^{(0)}$, $n^{(1)}$, ihre Brennweiten $b^{(0)}$, $b^{(1)}$, und die
Brennweite der gesuchten Linse $= x$. Wenn dann $q$ die totale
Brennweite des Fernrohrs bezeichnet, $s$ die Entfernung der ersten
Linse vom Gegenstande, $s'$ die Entfernung der letzten Linse vom
Bilde, und $d^{(0)}$, $d^{(1)}$ die Zwischendistanzen, so hat man die
Bedingungsgleichungen:
\[
\frac{1}{b^{(0)}} + \frac{1}{b^{(1)}} + \frac{1}{x}
- \frac{d^{(0)}}{b^{(0)}b^{(1)}} - \frac{d^{(1)}}{b^{(1)}x} =
\frac{1}{q},
\]
\[
s - \frac{ss'}{q} + s' = d^{(0)} + d^{(1)},
\]
und der Farbenfehler wird proportional
\[
\frac{1}{(n^{(0)}-1)\,b^{(0)}} + \frac{1}{(n^{(1)}-1)\,b^{(1)}}
+ \frac{1}{(n-1)\,x}.
\]
F\"ur ein Minimum unter der Nebenbedingung der gegebenen Gesamtl\"ange
leitet man durch Variation von $x$ ab und setzt die Variation des
Ausdrucks $= 0$.

\textbf{Zum \S 3.} \ L\"osung der Aufgabe mittels der Formel:
\[
\frac{1}{(n^{(0)}-1)\,b^{(0)2}}
+ \frac{1}{(n^{(1)}-1)\,b^{(1)2}}
+ \frac{1}{(n-1)\,x^{2}} = 0,
\]
wobei sich zeigt, da\ss\ die Aufl\"osung nur m\"oglich ist, wenn die
beiden gegebenen Gl\"aser von verschiedener Art sind, d.~h.~eines
Kronglas, das andere Flintglas.

\textbf{Zum \S 4.} \ Allgemeine Formel f\"ur die chromatische
Aberration:
\[
\partial q = q^{2}\left\{
\frac{\partial n^{(0)}}{(n^{(0)}-1)^{2}\,b^{(0)}}
+ \frac{\partial n^{(1)}}{(n^{(1)}-1)^{2}\,b^{(1)}}
+ \cdots
\right\}.
\]

\textbf{Zum \S 7.} \ Wenn $b^{(0)}$ und $b^{(1)}$ beliebig sind, so
nimmt der Farbenfehler seinen kleinsten Werth an f\"ur
\[
x = \sqrt{-\frac{(n^{(0)}-1)\,b^{(0)2} + (n^{(1)}-1)\,b^{(1)2}}{n-1}}.
\]

\textbf{Zum \S 8.} \ Bei einem doppelten Fernrohre, wo ein einfaches
Objektiv und ein einfaches Oeular gegeben sind, wird die L\"osung:
\[
\frac{1}{x} = -\frac{n^{(0)}-1}{n-1}\cdot\frac{1}{b^{(0)}}
- \frac{n^{(1)}-1}{n-1}\cdot\frac{1}{b^{(1)}}.
\]


\newpage
\section*{Erl\"auterungen zu den Abschnitten [I.] bis [III.]}
\addcontentsline{toc}{section}{Erl\"auterungen zu den Abschnitten [I.] bis [III.]}

% =====================================================
% ERLAUTERUNGEN ZU [I.]
% =====================================================
\subsection*{Erl\"auterung zu [I. Neuer Algorithmus]}

Der in [I.] entwickelte Algorithmus hat seinen Ursprung in der
Theorie der Kettenbr\"uche. Bezeichnet man n\"amlich den Kettenbruch
\[
a + \cfrac{1}{b + \cfrac{1}{c + \cfrac{1}{d + \cfrac{1}{\ddots}}}}
\]
durch $[a, b, c, d, \ldots]$, so ist bekanntlich
\[
[a, b, c, \ldots, z] = \frac{(a,z)}{(b,z)},
\]
wo die Gr\"ossen $(a,z)$ durch die im Text gegebene Recursionsformel
definiert sind. Die dortigen S\"atze I. und II. sind dann
unmittelbare Folgen der Theorie der convergentes.

\medskip

\textbf{Zu Formel [III.]--[VI].} \ Die dort eingef\"uhrten Gr\"ossen
$\alpha$, $\beta$, $\gamma$, $\delta$ mit ihren Accenten sind die
Coefficienten der linearen Transformation, welche die optische
Abbildung durch ein System von Linsen darstellt. Schreibt man die
gesammte Brechkraft des Systems
\[
\frac{1}{f} + \frac{1}{f'} + \frac{1}{f''} + \cdots
- \frac{h}{ff'} - \frac{h'}{f'f''} - \cdots
+ \frac{hh'}{ff'f''} + \cdots,
\]
so sind die Gr\"ossen $\alpha^{(\mu)}$, $\beta^{(\mu)}$,
$\gamma^{(\mu)}$, $\delta^{(\mu)}$ nichts anderes als die Elemente
derjenigen substitutions:\,formel, welche den Uebergang von den
Coordinaten des einfallenden Strahls zu denen des austretenden
vermittelt. Man hat n\"amlich, wenn $y$, $\eta$ die Coordinate und
Richtung des einfallenden, $y^{(\mu)}$, $\eta^{(\mu)}$ die des
austretenden Strahls bedeuten:
\[
y^{(\mu)} = \alpha^{(\mu)} y - \beta^{(\mu)} \eta,\qquad
\eta^{(\mu)} = -\gamma^{(\mu)} y + \delta^{(\mu)} \eta.
\]

\medskip

\textbf{Zu Formel [VII].} \ Die Formel f\"ur die Breite des farbigen
Saumes ergiebt sich durch folgende Betrachtung. Wenn $n$ das
Brechungsverh\"altniss f\"ur eine bestimmte Art von Licht, $n + \partial n$
dasjenige f\"ur eine benachbarte Art bezeichnet, so leitet man aus
den Hauptformeln der Dioptrik durch Differentiation nach $n$ die
Aenderung der Bildlage und der Vergr\"osserung ab. Die Differenz der
Bildlagen f\"ur die beiden Lichtarten gibt die Breite des farbigen
Randes. Der Ausdruck
\[
\alpha^{\mu}\,\partial\delta^{\mu} - \beta^{\mu}\,\partial\gamma^{\mu}
\]
ist die Summe derjenigen Gr\"ossen, welche von den einzelnen Linsen
her\"uhren, und kann durch die im Text angegebenen drei Formen
[VIIa], [VIIIb], [VIIc] dargestellt werden.

Die Formel [VIIa] erh\"alt man, indem man die Werthe von
$\partial\delta^{\mu}$ und $\partial\gamma^{\mu}$ durch die
Brechkrafte ausdr\"uckt und die Reductionsformeln des Algorithmus
anwendet. Die Formel [VIIIb] folgt durch partielle Umformung unter
Benutzung der Relation
\[
\alpha^{(k)}\delta^{(k)} - \beta^{(k)}\gamma^{(k)} = 1,
\]
welche f\"ur jedes $k$ identisch besteht. Die Formel [VIIc] endlich
entsteht durch eine analoge Umformung unter Vertauschung der Rollen
von $\alpha$, $\delta$ und $\beta$, $\gamma$.

% =====================================================
% ERLAUTERUNGEN ZU [II.]
% =====================================================
\subsection*{Erl\"auterung zu [II. Analyse von Fernr\"ohren]}

Die in diesem Abschnitte enthaltenen Formeln sind auf die
practische Construction von Fernr\"ohren berechnet. Die Formel f\"ur
die \textit{Feinheit der Theilung} bezieht sich auf die von
{\sc Fraunhofer} eingef\"uhrte Methode, die Gitter f\"ur
Theilungsstriche auf optischem Wege zu pr\"ufen.

Die Formel f\"ur die \textit{Unsch\"arfe}
\[
\tfrac{1}{4}\delta^{\mu}g\cdot\partial\frac{\beta^{\mu}}{\delta^{\mu}}
\]
giebt die Aenderung des Orts des Bildes in Bezug auf die Oeffnung,
oder die Aenderung von $\frac{\beta^{\mu}}{\delta^{\mu}}$ in R\"ucksicht
auf die Brechung. Der Factor $\delta^{\mu}$ bedeutet die
Vergr\"osserung, $g$ das Gesichtsfeld.

\textbf{Zu \S 1.} \ F\"ur drei Linsen mit gegebenen Distanzen und
gesuchten Brennweiten wird die L\"osung dadurch bestimmt, da\ss\ die
Oeffnungen der drei Linsen einander gleich werden sollen. Dies giebt
\[
y^{(1)} = y^{(2)} = y^{(3)} = \sqrt[3]{y^{2}f},
\]
wo $y$ die gesammte Oeffnung, $f$ die Brennweite des Objectivs ist.

\textbf{Zu \S 2.} \ Bei der Construction des dreifachen Fernrohrs
mit divergirender Mittellinse werden die Verh\"altnisse so bestimmt,
da\ss\ die sph\"arische Aberration und der chromatische Fehler
zugleich verschwinden. Die Formeln f\"ur $y^{(2)}$, $h'$,
$y^{(3)}$ sind daraus abgeleitet, da\ss\ die Aberrationen der drei
Linsen einzeln verschwinden und ihre Summe gleich Null wird.

% =====================================================
% ERLAUTERUNGEN ZU [III.]
% =====================================================
\subsection*{Erl\"auterung zu [III. Dioptrik]}

Die in diesem Abschnitte behandelten achromatischen und
aplanatischen Objective sind nach den von {\sc Fraunhofer}
eingef\"uhrten Principien berechnet. Die Zahlen f\"ur Kronglas und
Flintglas sind die von {\sc Fraunhofer} bestimmten.

\textbf{Zu \S 1--2.} \ Die Formel
\[
\frac{1}{f^{(1)}} = \frac{1}{v} + \frac{1}{t} + 1
\]
bezieht sich auf die Correction der chromatischen Aberration durch
eine Combination von drei Linsen. Die Gr\"ossen $v$ und $t$ sind die
reciprocen Brennweiten der beiden \oe usseren Linsen, $f^{(1)}$ die
der mittleren.

\textbf{Zu \S 8--9.} \ Die vier Gleichungen f\"ur das aplanatische
und achromatische System zweier Linsen sind:
\begin{enumerate}
\item Die Gleichung der chromatischen Correction:
\[
\frac{1}{(n-1)b} + \frac{1}{(n'-1)b'} = 0.
\]
\item Die Gleichung der sph\"arischen Aberration f\"ur die erste
Oberfl\"ache:
\[
\frac{n+2}{n}\cdot\frac{1}{r^{3}}
+ \frac{n'+2}{n'}\cdot\frac{1}{r'^{3}} = 0.
\]
\item Die Gleichung der sph\"arischen Aberration f\"ur die zweite
Oberfl\"ache.
\item Die Bedingung des verschwindenden Astigmatismus.
\end{enumerate}
Hieraus geht hervor, da\ss\ im Allgemeinen keine reelle L\"osung
existirt, und da\ss\ die M\"oglichkeit von der Wahl der Glassorten
abh\"angt.


\newpage
% =====================================================
% [V.] ACHROMATISCHE DOPPELAOBJECTIVE
% =====================================================
\subsection*{[V.]}
\begin{center}
\textbf{ACHROMATISCHE DOPPELAOBJECTIVE.}
\end{center}
\begin{center}
\rule{6cm}{0.4pt}\\[4pt]
[Aus Handbuch 21, das \"ubrige S.~83.]
\end{center}

Bei einem achromatischen Doppelaobjective werden zwei
biconvexe Linsen aus verschiedenem Glase so verbunden, da\ss\ sowohl
die chromatische Aberration als die sph\"arische ganz verschwinden.

Es seien f\"ur die erste Linse die Radien der Oberfl\"achen $r$ und
$r'$, f\"ur die zweite $\varrho$ und $\varrho'$; die Brechungsverh\"altnisse
seien $n$ und $n'$. Dann hat man zun\"achst die beiden Gleichungen
f\"ur die chromatische Correction:
\[
\text{(1)}\qquad \frac{1}{(n-1)b} + \frac{1}{(n'-1)b'} = \frac{1}{q},
\]
wo $q$ die gesammte Brennweite, und
\[
\text{(2)}\qquad \frac{1}{b} = (n-1)\left(\frac{1}{r} - \frac{1}{r'}\right),\qquad
\frac{1}{b'} = (n'-1)\left(\frac{1}{\varrho} - \frac{1}{\varrho'}\right).
\]

Die sph\"arische Aberration wird f\"ur jede Linse durch die bekannte
Formel ausgedr\"uckt:
\[
\text{(3)}\qquad
\frac{n+2}{n(n-1)^{2}}\left(\frac{1}{r} - \frac{n+1}{n+2}\cdot\frac{1}{b}\right)^{2}
+ \frac{1}{n}\cdot\frac{1}{r^{2}} = K,
\]
wo $K$ eine Constante ist, die f\"ur die erste Linse einen Werth,
f\"ur die zweite den entgegengesetzten haben mu\ss, damit die
Gesamtaberration verschwinde.

\medskip

\textbf{Satz.} \ Wenn die beiden Linsen sich ber\"uhren, so ist die
Bedingung der vollst\"andigen Aplanasie
\[
\frac{n^{2}(n'-1)^{2}}{n'^{2}(n-1)^{2}} =
\frac{n+2}{n'+2}\cdot\frac{n'^{2}-1}{n^{2}-1}.
\]

\textbf{Beweis.} \ Setzt man in (3) $r' = \infty$, so wird die
Aberration der ersten Linse
\[
K = \frac{n+2}{n(n-1)^{2}}\cdot\frac{1}{r^{2}}
- 2\frac{n+2}{n(n-1)^{2}}\cdot\frac{n+1}{n+2}\cdot\frac{1}{rb}
+ \left\{\frac{(n+1)^{2}}{n(n-1)^{2}} + \frac{1}{n}\right\}\frac{1}{b^{2}}.
\]
Durch Umformung erh\"alt man:
\[
K = \frac{n+2}{n(n-1)^{2}}\left(\frac{1}{r} - \frac{n+1}{n+2}\cdot\frac{1}{b}\right)^{2}
+ \frac{1}{n(n-1)^{2}}\cdot\frac{1}{b^{2}}.
\]
F\"ur die zweite Linse ergiebt sich ein analoger Ausdruck $K'$, und
die Bedingung $K + K' = 0$ f\"uhrt auf die im Satze gegebene
Gleichung.

\medskip

Setzt man die numerischen Werthe f\"ur Kronglas ($n = 1.53$) und
Flintglas ($n' = 1.58$), so ergiebt sich, da\ss\ die Bedingung nicht
genau erf\"ullt ist. Jedoch kann man durch eine kleine Ab\"anderung
der Brechungsverh\"altnisse oder durch Einf\"uhrung eines
zwischenliegenden Mediums die Gleichung befriedigen.

\medskip

\textbf{Zum \S 5.} \ Die Gleichungen (1) und (2) geben:
\[
\frac{1}{r} - \frac{1}{r'} = \frac{1}{(n-1)b},\qquad
\frac{1}{\varrho} - \frac{1}{\varrho'} = \frac{1}{(n'-1)b'}.
\]
Wenn die Linsen sich ber\"uhren, ist $d = 0$, also $b + b' = q$.

\textbf{Zum \S 6.} \ Die Berechnung der sph\"arischen Aberration f\"ur
eine einzelne Linse mit den Radien $r$, $r'$ giebt die Formel:
\[
\text{Aberration} =
y^{4}\left\{\frac{n+2}{4n(n-1)^{2}}\left(\frac{1}{r} - \frac{n+1}{n+2}\cdot\frac{1}{b}\right)^{2}
+ \frac{1}{4n(n-1)^{2}}\cdot\frac{1}{b^{2}}\right\},
\]
wo $y$ die halbe Oeffnung ist.

\textbf{Zum \S 8.} \ Bei Ber\"ucksichtigung der Dicke der Linsen
werden die Formeln complicirter. Man hat dann statt der einfachen
Bedingung $K + K' = 0$ eine Gleichung von der Form:
\[
K + K' + \frac{d}{b^{2}}\left(\frac{1}{b} - \frac{1}{b'}\right) = 0,
\]
wo $d$ die Dicke der ersten Linse ist.

\textbf{Zum \S 10.} \ Nimmt man die von {\sc Fraunhofer} bestimmten
Werthe:
\[
n = 1.5258,\quad n' = 1.5852,
\]
so wird
\[
\frac{n^{2}(n'-1)^{2}}{n'^{2}(n-1)^{2}} = 1.073,\qquad
\frac{n+2}{n'+2}\cdot\frac{n'^{2}-1}{n^{2}-1} = 1.064.
\]
Die Differenz ist etwa $0.009$, also von der Ordnung $\frac{1}{100}$.
Dieser Unterschied erkl\"art, warum die vollkommene Achromasie und
Aplanasie nicht gleichzeitig erreicht werden kann, ohne die
Glassorte zu ver\"andern.


% =====================================================
% ERLAUTERUNGEN ZU [V.]
% =====================================================
\subsection*{Erl\"auterung zu [V. Achromatische Doppelaobjective]}

Die im Text entwickelte Theorie der achromatischen
Doppelaobjective ist von fundamentaler Bedeutung f\"ur die
Construction von Fernr\"ohren und Mikroskopen. Die dortige Formel
\[
\frac{n^{2}(n'-1)^{2}}{n'^{2}(n-1)^{2}} =
\frac{n+2}{n'+2}\cdot\frac{n'^{2}-1}{n^{2}-1}
\]
enth\"alt die nothwendige Bedingung daf\"ur, da\ss\ zwei sich
ber\"uhrende Linsen aus verschiedenem Glase zugleich achromatisch
und aplanatisch sein k\"onnen.

\textbf{Zu den Formeln (1)--(3).} \ Die Gleichung (1) dr\"uckt die
bekannte Bedingung aus, da\ss\ die algebraische Summe der
Dispersionen verschwinde. Die Gleichungen (2) sind die
Grundgleichungen der Dioptrik f\"ur einzelne Linsen. Die Formel (3)
endlich ist die von {\sc Gauss} gegebene Ausdrucksform f\"ur die
sph\"arische Aberration einer Linse.

\textbf{Zum Beweis.} \ Die Formel (3) leitet sich aus der Entwicklung
der Brechungsformel nach Potenzen der Oeffnung ab. Setzt man
\[
\sin\varphi = \varphi - \frac{\varphi^{3}}{6} + \cdots,
\]
so erh\"alt man f\"ur den Schnittpunkt der gebrochenen Strahlen mit
der Axe einen Ausdruck von der Form
\[
s = s_{0} + Ay^{2} + By^{4} + \cdots,
\]
wo $s_{0}$ der Gau\ss'sche Bildpunkt, $y$ die halbe Oeffnung ist.
Der Coefficient $A$ ist die sph\"arische Aberration erster Ordnung,
und die Formel (3) giebt den genauen Werth dieses Coefficienten
unter Ber\"ucksichtigung beider Fl\"achen der Linse.

\textbf{Zu den Zahlenwerthen.} \ Die numerische Vergleichung zeigt,
da\ss\ f\"ur die gew\"ohnlichen Glassorten (Kronglas und Flintglas)
die Bedingung der vollkommenen Aplanasie nicht erf\"ullt ist. Dies
hat zur Folge, da\ss\ man in der Praxis entweder einen kleinen
Farbfehler zul\"a\ss t, oder durch Einf\"uhrung einer dritten Linse
(eines Mittelgliedes) die vollkommene Correction erstrebt.

\textbf{Zu \S 8.} \ Die Formel f\"ur die Ber\"ucksichtigung der Dicke
leitet sich aus der allgemeinen Formel f\"ur die sph\"arische
Aberration einer dicken Linse ab. Wenn $d$ die Dicke und $n$ das
Brechungsverh\"altniss ist, so hat man statt der einfachen Formel
die genauere:
\[
\frac{1}{b} = (n-1)\left(\frac{1}{r_{1}} - \frac{1}{r_{2}}
+ \frac{n-1}{n}\cdot\frac{d}{r_{1}r_{2}}\right).
\]
Dies giebt die Correction f\"ur die Dicke in der Formel f\"ur die
Brennweite, und analoge Correctionen treten in der Aberrationsformel
auf.


% =====================================================
% [VI.] REFLEXION VON DER HINTEREN FLAECHE EINER LINSE
% =====================================================
\subsection*{[VI.]}
\begin{center}
\textbf{REFLEXION VON DER HINTEREN FL\"ACHE EINER LINSE.}
\end{center}
\begin{center}
\rule{6cm}{0.4pt}\\[4pt]
[Aus Handbuch 21, S.~84.]
\end{center}

Wenn Lichtstrahlen durch eine Linse gehen und von der hinteren
gebogenen Fl\"ache theilweise zur\"uckgeworfen werden, so entstehen
besondere Erscheinungen, die f\"ur die Theorie der optischen
Instrumente von Wichtigkeit sind.

Es seien f\"ur eine einzelne biconvexe Linse die Radien $r$ und $r'$,
das Brechungsverh\"altniss $n$, die Dicke $d$. Ein von der
vorderen Fl\"ache eintretender Strahl werde an der hinteren Fl\"ache
reflectirt und trete wieder durch die vordere Fl\"ache aus. Der
Ort des entstehenden Spiegelbildes l\"a\ss t sich durch folgende
Ueberlegung bestimmen.

\textbf{Satz.} \ Das von der hinteren Fl\"ache einer Linse erzeugte
Spiegelbild hat die Entfernung von der vorderen Fl\"ache:
\[
s' = \frac{r'}{2n}\cdot\frac{2n-1}{n-1} - \frac{d}{n},
\]
wo $r'$ der Radius der hinteren Fl\"ache ist.

\textbf{Beweis.} \ Ein von der vorderen Fl\"ache nach der hinteren
laufender Strahl habe die Richtung $\eta$. Nach der Brechung an der
ersten Fl\"ache wird er die hintere Fl\"ache unter dem Winkel
\[
\sin\psi = n\sin\varphi
\]
treffen, wo $\varphi$ der Einfallswinkel an der ersten Fl\"ache ist.
Nach der Reflexion an der zweiten Fl\"ache kehrt er zur\"uck und wird
an der ersten Fl\"ache wieder gebrochen. Die gesammte Wegl\"ange
innerhalb der Linse ist $2d/\cos\psi'$, wo $\psi'$ der Winkel im
Glase ist. Durch Entwicklung nach Potenzen der Oeffnung erh\"alt man
f\"ur den Schnittpunkt mit der Axe den oben angegebenen Werth.

\medskip

\textbf{Zu \S 2.} \ Wenn die hintere Fl\"ache versilbert ist, so
wirkt die Linse als Spiegel mit einer bestimmten Brennweite, die
sich aus der obigen Formel ergiebt. Diese Anordnung wird bei
gewissen Spiegelteleskopen angewendet.

\textbf{Zu \S 3.} \ Die Intensit\"at des reflectirten Lichtes h\"angt
von den Reflexionsverm\"ogen der Fl\"ache ab. Nach den
{Fresnel}schen Formeln ist f\"ur senkrechten Einfall das
Reflexionsverm\"ogen
\[
R = \left(\frac{n-1}{n+1}\right)^{2}.
\]
F\"ur Kronglas ($n = 1.53$) ergiebt sich $R = 0.043$, d.~h.~etwa
$4\,\frac{1}{3}\%$ des einfallenden Lichtes wird zur\"uckgeworfen.

\textbf{Zu \S 5.} \ Bei einer Combination von Linsen, wie sie in
Objectiven vorkommt, wird das von jeder hinteren Fl\"ache
reflectirte Licht ein mehr oder weniger scharfes Bild der
Objektebene erzeugen. Diese Erscheinung ist bekannt unter dem Namen
\textit{reflexe Bilder} oder \textit{Nebenbilder} und st\"ort bei
photographischen Aufnahmen sowie bei genauen Messungen.


% =====================================================
% ERLAUTERUNGEN ZU [VI.]
% =====================================================
\subsection*{Erl\"auterung zu [VI. Reflexion von der hinteren Fl\"ache einer Linse]}

Die in diesem Abschnitte behandelte Erscheinung tritt bei allen
Linsensystemen auf, wo die Fl\"achen nicht entspiegelt sind. Die
Formel f\"ur die Lage des Spiegelbildes leitet sich aus der
Verfolgung eines einzelnen Strahls durch die Linse und seiner
Reflexion an der hinteren Fl\"ache ab.

\textbf{Zu den Formeln.} \ Die Formel
\[
s' = \frac{r'}{2n}\cdot\frac{2n-1}{n-1} - \frac{d}{n}
\]
gilt f\"ur den Fall, da\ss\ das Object im Unendlichen liegt und die
Linse biconvex ist. F\"ur ein endliches Object mit der Entfernung
$s$ von der vorderen Fl\"ache wird die Bildweite:
\[
s'' = \frac{r'}{2}\left(\frac{1}{n} - \frac{1}{2n-1}\right)
- \frac{d}{2n-1} + \frac{s}{2n-1}.
\]

\textbf{Zu \S 3.} \ Die {\sc Fresnel}schen Formeln f\"ur das
Reflexionsverm\"ogen bei beliebigem Einfallswinkel $\varphi$ sind:
\[
R_{s} = \left(\frac{\sin(\varphi-\psi)}{\sin(\varphi+\psi)}\right)^{2},\qquad
R_{p} = \left(\frac{\tan(\varphi-\psi)}{\tan(\varphi+\psi)}\right)^{2},
\]
wo $\psi$ der Brechungswinkel ist. F\"ur senkrechten Einfall
vereinfachen sich diese zu der im Text gegebenen Formel.


\newpage
\section*{Briefwechsel}
\addcontentsline{toc}{section}{Briefwechsel}

% =====================================================
% [1.] SCHUMACHER AN GAUSS
% =====================================================
\subsection*{[1.]}
\begin{center}
\textbf{Schumacher an Gau\ss.} Altona, 11.~September 1810.
\end{center}
\begin{center}
\rule{6cm}{0.4pt}\\[4pt]
[Briefwechsel zwischen Gau\ss\ und Schumacher I, S.~534.]
\end{center}

Sie schreiben mir in Ihrem letzten Briefe: \glqq Ihre Anzeige,
die ich in einem gewissen Aufsatze erw\"ahnt gefunden habe, erwartet
wohl kaum einer von mir.\grqq\ Das bezog sich auf eine fr\"uhere
\"ausserung meinerseits \"uber die Anwendung der von mir in der
{Theoria Motus} gegebenen Formeln auf die Berechnung der Cometen- und
Planetenbahnen. Ich hatte mir erlaubt zu bemerken, da\ss\ diese
Formeln, obgleich in theoretischer Hinsicht sehr elegant, f\"ur die
practische Anwendung nicht immer bequem seien, und da\ss\ ich daher
andere Methoden vorziehe.

Sie erwidern darauf: \glqq Die Methode, welche Sie vorziehen, ist
ohne Zweifel die bekannte, auf den drei beobachteten Oertern
beruhende. Ich gebe zu, da\ss\ diese Methode in vielen F\"allen
bequemer ist, allein sie hat den Nachtheil, da\ss\ sie nicht allgemein
anwendbar ist, namentlich nicht, wenn die drei Oerter zu nahe bei
einander liegen. Meine Methode dagegen ist ganz allgemein und
leidet unter dieser Beschr\"ankung nicht.\grqq

Hierzu bemerke ich, da\ss\ der Einwand gegen die gew\"ohnliche Methode
wohl begr\"undet ist, wenn die drei Oerter sehr nahe beieinander
liegen, da\ss\ aber in den meisten praktischen F\"allen diese
Bedingung nicht erf\"ullt ist, und da\ss\ daher die gew\"ohnliche
Methode f\"ur die Anwendung meist ausreicht.

\medskip

\textbf{Anmerkung.} \ Der in diesem Briefe erw\"ahnte Aufsatz von
{\sc Schumacher} betrifft wohl die im \textit{Monatlichen
Correspondenz} erschienenen Untersuchungen \"uber die
Bahnbestimmung der Cometen.


% =====================================================
% [2.] GAUSS AN SCHUMACHER
% =====================================================
\subsection*{[2.]}
\begin{center}
\textbf{Gau\ss\ an Schumacher.} G\"ottingen, 18.~October 1810.
\end{center}
\begin{center}
\rule{6cm}{0.4pt}\\[4pt]
[Briefwechsel zwischen Gau\ss\ und Schumacher I, S.~536.]
\end{center}

Ich danke Ihnen f\"ur Ihren freundschaftlichen Brief. Was die
Anwendung meiner Formeln betrifft, so bin ich weit davon entfernt,
sie als die allein zweckm\"a\ss ige Methode f\"ur alle F\"alle
hinzustellen. Ich habe sie selbst nur f\"ur solche F\"alle
empfohlen, wo die gew\"ohnlichen Methoden versagen oder zu
unsicher sind. F\"ur die gew\"ohnliche Bahnbestimmung aus drei
Oertern halte ich die bekannten Methoden von {\sc Olbers} und
{\sc Eschenbach} f\"ur v\"ollig ausreichend.

Was ich jedoch f\"ur wichtig halte, ist die theoretische
Einsicht in die Natur der Aufgabe, welche sich durch meine Formeln
besser gewinnen l\"a\ss t als durch die gew\"ohnlichen. Die Formeln
zeigen n\"amlich deutlich, wie die Aufl\"osung von der L\"osung einer
Gleichung h\"oheren Grades abh\"angt, und wie sich diese Gleichung
in besonderen F\"allen vereinfacht.

\medskip

\textbf{Anmerkung.} \ Die in diesem Briefe erw\"ahnten Methoden von
{\sc Olbers} und {\sc Eschenbach} sind die bekannten
Bahnbestimmungsmethoden, welche auf der Annahme eines parabolischen
oder elliptischen Bogenelementes beruhen.


% =====================================================
% [3.] OLBERS AN GAUSS
% =====================================================
\subsection*{[3.]}
\begin{center}
\textbf{Olbers an Gau\ss.} Bremen, 5.~November 1810.
\end{center}
\begin{center}
\rule{6cm}{0.4pt}\\[4pt]
[{\it Wilhelm Olbers, Sein Leben und seine Werke} II, 1, S.~308.]
\end{center}

Ich habe mit grossem Interesse die Discussion zwischen Ihnen und
{\sc Schumacher} \"uber die Bahnbestimmungsmethoden verfolgt. Sie
werden sich erinnern, da\ss\ ich schon fr\"uher die Bemerkung
gemacht habe, da\ss\ Ihre Methode zwar sehr elegant, aber f\"ur den
praktischen Astronomen nicht immer bequem sei.

Was mich jedoch in diesem Zusammenhange besonders interessirt, ist
die Frage nach der Genauigkeit der verschiedenen Methoden. Sie
wissen, da\ss\ ich bei meiner Methode der Bahnbestimmung stets die
Mittel anwende, die Genauigkeit durch successive Approximation zu
vergr\"ossern. Wie steht es aber mit der analytischen Bestimmung
der Fehlergrenzen?

Ich habe neulich einen Aufsatz von {\sc Laplace} \"uber diesen
Gegenstand gelesen, worin er zeigt, da\ss\ die Methode der kleinsten
Quadrate die besten Resultate giebt. Halten Sie es f\"ur m\"oglich,
diese Methode auf die Bahnbestimmung anzuwenden?

\medskip

\textbf{Anmerkung.} \ Die von {\sc Olbers} erw\"ahnte Arbeit von
{\sc Laplace} ist wohl die im \textit{Th\'eorie analytique des
probabilit\'es} enthaltene Theorie der Fehlervertheilung.


% =====================================================
% [4.] GAUSS AN OLBERS
% =====================================================
\subsection*{[4.]}
\begin{center}
\textbf{Gau\ss\ an Olbers.} G\"ottingen, 20.~November 1810.
\end{center}
\begin{center}
\rule{6cm}{0.4pt}\\[4pt]
[{\it Wilhelm Olbers, Sein Leben und seine Werke} II, 1, S.~310.]
\end{center}

Ihre Frage nach der Anwendung der Methode der kleinsten Quadrate
auf die Bahnbestimmung ist sehr wichtig und ich habe dar\"uber
schon lange nachgedacht. Die Schwierigkeit besteht darin, da\ss\ die
Beziehungen zwischen den beobachteten Gr\"ossen und den gesuchten
Elementen nicht linear sind, und da\ss\ daher die directe Anwendung
der Methode der kleinsten Quadrate auf sehr complicirte Gleichungen
f\"uhrt.

Ich habe jedoch einen Weg gefunden, diese Schwierigkeit zu umgehen.
Man kann n\"amlich die Aufgabe in zwei Theile zerlegen: zuerst
bestimmt man die Elemente n\"aherungsweise nach einer der bekannten
Methoden, und dann verbessert man sie durch Anwendung der Methode
der kleinsten Quadrate auf die Abweichungen der beobachteten Oerter
von den berechneten. Diese Methode habe ich in meiner
\textit{Theoria Motus} ausf\"uhrlich entwickelt.

\medskip

\textbf{Anmerkung.} \ Die hier beschriebene Methode ist die in der
\textit{Theoria Motus} (1809) gegebene Theorie der Bahnbestimmung,
wo zuerst eine n\"aherungsweise L\"osung gefunden und dann durch
successive Approximation verbessert wird.


% =====================================================
% [5.] SCHUMACHER AN GAUSS
% =====================================================
\subsection*{[5.]}
\begin{center}
\textbf{Schumacher an Gau\ss.} Altona, 10.~December 1810.
\end{center}
\begin{center}
\rule{6cm}{0.4pt}\\[4pt]
[Briefwechsel zwischen Gau\ss\ und Schumacher I, S.~540.]
\end{center}

Ich bin sehr erfreut zu h\"oren, da\ss\ Sie die Anwendung der
Methode der kleinsten Quadrate auf die Bahnbestimmung so erfolgreich
durchgef\"uhrt haben. Ich habe selbst versucht, nach Ihrer Methode
die Bahn des letzten Cometen zu berechnen, und bin zu sehr guten
Resultaten gekommen.

Es giebt jedoch Einen Punkt, \"uber den ich noch Ihre Meinung h\"oren
m\"ochte. Sie nehmen bei Ihrer Methode an, da\ss\ die
n\"aherungsweise bestimmten Elemente so genau sind, da\ss\ die
Abweichungen linear genommen werden k\"onnen. Was aber, wenn dies
nicht der Fall ist? M\"ussen wir dann nicht auf die urspr\"ungliche
nichtlineare Gleichung zur\"uckgehen?

\medskip

\textbf{Anmerkung.} \ Die Frage von {\sc Schumacher} betrifft die
Convergenz der successiven Approximationen bei der Bahnbestimmung.
{\sc Gauss} hat sp\"ater gezeigt, da\ss\ die Linearisierung unter
sehr allgemeinen Voraussetzungen zul\"assig ist.


% =====================================================
% [6.] GAUSS AN SCHUMACHER
% =====================================================
\subsection*{[6.]}
\begin{center}
\textbf{Gau\ss\ an Schumacher.} G\"ottingen, 15.~Januar 1811.
\end{center}
\begin{center}
\rule{6cm}{0.4pt}\\[4pt]
[Briefwechsel zwischen Gau\ss\ und Schumacher I, S.~542.]
\end{center}

Ihre Bedenklichkeit gegen die Linearisierung ist theoretisch
wohl begr\"undet, aber in der Praxis kaum von Bedeutung. Wenn die
n\"aherungsweise bestimmten Elemente nur einigerma\ss en genau sind
--- und dies wird bei den bekannten Methoden stets der Fall sein,
sobald man nicht gerade die allerschlechtesten Beobachtungen hat ---
so convergirt die successive Approximation sehr schnell, und man
braucht h\"ochstens zwei oder drei Correctionsschritte.

Ich habe \"ubrigens bemerkt, da\ss\ man die Convergenz noch
beschleunigen kann, indem man statt der einfachen Correction eine
verbesserte Formel anwendet, welche die Glieder zweiter Ordnung
mitber\"ucksichtigt. Die Entwicklung dieser Formel w\"urde jedoch zu
weit f\"uhren, und ich behalte sie mir f\"ur eine sp\"atere
Gelegenheit vor.

\medskip

\textbf{Anmerkung.} \ Die hier erw\"ahnte verbesserte Formel ist
wahrscheinlich die sp\"ater in der \textit{Theoria Motus} Art.~172
ff.\ gegebene Entwicklung der Correctionen zweiter Ordnung.


% =====================================================
% [7.] OLBERS AN GAUSS
% =====================================================
\subsection*{[7.]}
\begin{center}
\textbf{Olbers an Gau\ss.} Bremen, 1.~Februar 1811.
\end{center}
\begin{center}
\rule{6cm}{0.4pt}\\[4pt]
[{\it Wilhelm Olbers, Sein Leben und seine Werke} II, 1, S.~312.]
\end{center}

Ich danke Ihnen herzlich f\"ur Ihre ausf\"uhrliche Erkl\"arung.
Ich bin nun ganz beruhigt und werde bei meinen k\"unftigen
Bahnbestimmungen Ihre Methode der successive Approximation mit der
Methode der kleinsten Quadrate anwenden.

Es giebt noch Einen Punkt, den ich gerne zur Sprache bringen
m\"ochte. Sie haben in Ihrer \textit{Theoria Motus} die Formeln
f\"ur die elliptische Bewegung sehr elegant entwickelt. Halten Sie
es f\"ur m\"oglich, analoge Formeln auch f\"ur die parabolische und
hyperbolische Bewegung zu geben?

Ich weiss, da\ss\ die Formeln f\"ur den parabolischen Fall schon
lange bekannt sind, aber f\"ur den hyperbolischen Fall scheint mir
die Sache noch nicht vollst\"andig erledigt zu sein.

\medskip

\textbf{Anmerkung.} \ Die von {\sc Olbers} erw\"ahnten Formeln f\"ur
die hyperbolische Bewegung sind sp\"ater von {\sc Gauss} in der
\textit{Theoria Motus} Art.~100 ff.\ vollst\"andig entwickelt
worden.


% =====================================================
% [8.] GAUSS AN OLBERS
% =====================================================
\subsection*{[8.]}
\begin{center}
\textbf{Gau\ss\ an Olbers.} G\"ottingen, 12.~Februar 1811.
\end{center}
\begin{center}
\rule{6cm}{0.4pt}\\[4pt]
[{\it Wilhelm Olbers, Sein Leben und seine Werke} II, 1, S.~314.]
\end{center}

Was die Formeln f\"ur die hyperbolische Bewegung betrifft, so
kann ich Ihnen versichern, da\ss\ ich dieselben vollst\"andig
entwickelt habe und da\ss\ sie nicht schwieriger sind als die f\"ur
die elliptische Bewegung. Der einzige Unterschied besteht darin,
da\ss\ statt der trigonometrischen Functionen die hyperbolischen
Functionen auftreten.

Man hat n\"amlich f\"ur die hyperbolische Bewegung an Stelle der
{Kepler}schen Gleichung:
\[
e\sinh F - F = \frac{k\sqrt{2}}{a^{3/2}}(t - T),
\]
wo $e$ die Excentricit\"at, $a$ die halbe Hauptaxe, $F$ die
analoge Excentricit\"atsanomalie und $T$ die Zeit des Periheldurchgangs
ist. Die Aufl\"osung dieser Gleichung geschieht ganz nach denselben
Principien wie bei der {Kepler}schen Gleichung, und ich habe
hierf\"ur eine sehr bequeme Methode gefunden.

Die Entwicklung der Formeln f\"ur die Coordination und ihre
Ableitungen nach der Zeit ergiebt sich dann durch ganz analoge
Schl\"usse wie im elliptischen Fall, und man erh\"alt schlie\ss lich
Formeln, welche in jeder Hinsicht mit den elliptischen
parallel laufen.

Ich werde diese Formeln in einem sp\"ateren Theile meiner
\textit{Theoria Motus} ver\"offentlichen.

\medskip

\textbf{Anmerkung.} \ Die hier erw\"ahnte Arbeit \"uber die
hyperbolische Bewegung ist in der \textit{Theoria Motus} Art.~100
ff.\ enthalten. Die dort gegebene Aufl\"osung der hyperbolischen
{Kepler}schen Gleichung ist eine der ersten Anwendungen der
hyperbolischen Functionen auf ein astronomisches Problem.

\end{document}
